\documentclass[../main.tex]{subfiles}
\begin{document}
\section{Introduction}
\label{sec:detection_intro}
Object detection represents a significant leap in complexity beyond simple image classification. Instead of assigning a single label to an entire image, object detection aims to solve a dual problem: identifying \textit{what} objects are present in an image and precisely localizing \textit{where} they are located using \textbf{bounding boxes}. This capability for fine-grained scene understanding is critical for applications that require interaction with the physical world, most notably autonomous driving. In this context, the reliable detection of traffic signs, pedestrians, and other vehicles is a fundamental prerequisite for safe navigation.

To align our research with this critical application, our experiments are centered on the \textbf{LISA (Laboratory for Intelligent \& Safe Automobiles) Traffic Sign Dataset}. This dataset is specifically designed to benchmark the performance of object detectors on real-world traffic signs captured under a variety of conditions, making it an ideal environment to assess the practical security risks posed by adversarial attacks.

For our model, we have selected \textbf{YOLOv8n (You Only Look Once, version 8 nano)}. As a one-stage detector, the YOLO architecture is renowned for its exceptional balance of speed and accuracy, making it a leading choice for real-time applications where low-latency inference is paramount. The 'nano' variant, in particular, is optimized for efficiency, making it highly suitable for deployment on resource-constrained hardware, such as the on-board computers in autonomous vehicles. Our trained YOLOv8n model establishes a strong performance baseline on the clean LISA dataset, achieving a mean Average Precision (\textbf{mAP\textsubscript{50-95}}) of \textbf{81.1\%} and a \textbf{mAP\textsubscript{50}} of \textbf{92.7\%}. This strong detection performance, coupled with a rapid inference time of approximately \textbf{3.9 milliseconds} per image on our hardware, confirms its suitability for real-world deployment.

Given that such a high-performing and efficient model is a prime candidate for use in safety-critical systems, this chapter will investigate its robustness against adversarial attacks. We will explore methods designed to fool object detectors into failing to detect signs or misclassifying them, and subsequently evaluate defenses intended to mitigate these critical vulnerabilities.\section{Attacks}
\subsection{Evaluation Methodology}
To assess the effectiveness of adversarial attacks, including the Fast Gradient Sign Method (FGSM), on the YOLO object detection model, a systematic evaluation framework was designed. The primary goal was to quantify the impact of perturbations on detection performance while ensuring the methodology is reproducible and aligned with standard practices in computer vision research.

The evaluation was conducted on a test dataset of traffic sign images, with a fixed confidence threshold of 0.25 to filter predictions. This threshold was selected to balance precision and recall, ensuring that only highly confident detections were considered while minimizing false positives. Performance was measured using the mean Average Precision (mAP) metric, with a focus on mAP@0.5, which evaluates detection accuracy at an Intersection over Union (IoU) threshold of 0.5. This metric is widely adopted in object detection benchmarks (e.g., Pascal VOC, COCO) and provides a practical measure of model robustness under moderate localization requirements. Additionally, mAP@0.5:0.95 was computed, averaging precision across IoU thresholds from 0.5 to 0.95 in steps of 0.05, to assess the attack's impact on both detection and precise localization, offering a more stringent evaluation where relevant.

Class-specific mAP@0.5 was evaluated to identify vulnerabilities across different traffic sign categories (e.g., Speed Limit signs, Stop signs), with the drop in mAP ($\Delta$ mAP) calculated to highlight the relative degradation for each class. This approach enables a detailed understanding of how visual saliency and feature distinctiveness affect robustness. The methodology was applied consistently across all attacks to ensure comparability, with results presented in subsequent sections to avoid redundancy in attack-specific discussions.
\subsection{White-Box Attacks}
\subsubsection{FGSM Attack on YOLO}

\paragraph{Introduction}\leavevmode\\
The Fast Gradient Sign Method (FGSM) is a single-step adversarial attack that generates adversarial examples by perturbing inputs along the gradient of the loss function \cite{goodfellow2015explaining}. For YOLO \cite{lu2017noneed}, FGSM targets high-confidence detections from raw model outputs before non-maximum suppression (NMS) to maximize classification errors, causing misclassifications or missed detections. The attack uses an $\ell_\infty$-bounded perturbation (controlled by $\epsilon$) to maintain similarity while disrupting the model's predictions.

\paragraph{Methodology}\leavevmode\\

\begin{breakablealgorithm}{Fast Gradient Sign Method (FGSM) Attack on YOLO}
\begin{algorithmic}[1]
\Require $x \in \mathbb{R}^{H \times W \times C}$: input image, $model$: YOLO model, $\epsilon > 0$: perturbation magnitude, $\tau > 0$: confidence threshold
\Ensure $x_{adv} \in \mathbb{R}^{H \times W \times C}$: adversarial image
\State Preprocess $x$ for input to $model$ \Comment{E.g., normalize to $[0, 1]$, resize}
\State Compute model predictions: $p = model(x)$ \Comment{$p$ includes bounding boxes, objectness scores, class probabilities}
\State Extract detections $\mathcal{D} = \{ (b_i, s_i, c_i) \mid s_i \cdot \max(c_i) > \tau \}$, where $b_i$ are bounding boxes, $s_i$ are objectness scores, $c_i \in \mathbb{R}^K$ are class probabilities for $K$ classes
\If{$\mathcal{D} = \emptyset$}
    \State \Return None \Comment{No high-confidence detections}
\EndIf
\State Select target class labels $y_i = \arg\max_k (c_i)_k$ for each detection $(b_i, s_i, c_i) \in \mathcal{D}$
\State Define classification loss: $L(x, \mathcal{D}) = \sum_{(b_i, s_i, c_i) \in \mathcal{D}} \ell(c_i, y_i)$, where $\ell$ is cross-entropy loss
\State Compute gradient: $g = \nabla_x L(x, \mathcal{D}) \in \mathbb{R}^{H \times W \times C}$ \Comment{Gradient of loss w.r.t. $x$}
\State Compute perturbation: $\eta = \epsilon \cdot \text{sign}(g)$ \Comment{Scaled sign of gradient}
\State Generate adversarial image: $x_{adv} = x + \eta$
\State Clamp $x_{adv}$ to valid range: $x_{adv} = \text{clip}(x_{adv}, 0, 1)$ \Comment{Ensure pixel values in $[0, 1]$}
\State \Return $x_{adv}$
\end{algorithmic}
\end{breakablealgorithm}
\paragraph{Results}\leavevmode\\

The FGSM attack was applied with perturbation magnitudes $\epsilon=0.01$, $0.05$, and $0.1$. Figure~\ref{fig:fgsm_example_yolo} illustrates the effect of the attack on a "Speed Limit 120" sign with $\epsilon=0.1$, showing the original image, the adversarial image, and the pixel-wise difference.

% Figure: FGSM Example Image
\begin{figure}[H]
    \centering
    \includegraphics[width=0.8\textwidth]{images/CV/OD/fgsm.png} % Replace with the actual path to your image file
    \caption{FGSM Example showing the original, adversarial, and difference images.}
    \label{fig:fgsm_example_yolo}
\end{figure}

  Although the relatively large epsilon ($\epsilon=0.1$), the adversarial image is not much different from the original image, as the perturbations are small and the original detections disappeared. 

The following tables ~\ref{tab:yolo_fgsm_overall_performance}and ~\ref{tab:yolo_fgsm_class_results}present the overall performance metrics and class-specific mean Average Precision (mAP@0.5), highlighting the impact of adversarial attacks on detection performance.



% Table 1: Overall Performance
\begin{table}[H]
\caption{Overall Performance of YOLO under Clean and FGSM Attacks}
    \label{tab:yolo_fgsm_overall_performance}
    \centering
    \footnotesize
    \begin{tabular}{@{\hspace{0.3cm}}lS[table-format=1.3]S[table-format=1.3]S[table-format=1.3]S[table-format=1.3]@{\hspace{0.3cm}}}
        \toprule
        \textbf{Metric} & {\textbf{Clean}} & {\textbf{FGSM $\epsilon=0.01$}} & {\textbf{FGSM $\epsilon=0.05$}} & {\textbf{FGSM $\epsilon=0.1$}} \\
        \midrule
        Precision     & 0.935 & 0.893 & 0.722 & 0.582 \\
        Recall        & 0.854 & 0.810 & 0.622 & 0.491 \\
        mAP@0.5       & 0.927 & 0.836 & 0.595 & 0.416 \\
        mAP@0.5:0.95  & 0.811 & 0.713 & 0.492 & 0.341 \\
        \bottomrule
    \end{tabular}
\end{table}

% Table 2: mAP@0.5 per Class
\begin{table}[H]
\caption{Class-Specific mAP@0.5 under Clean and FGSM Attacks}
    \label{tab:yolo_fgsm_class_results}
    \centering
    \footnotesize
    \begin{tabular}{@{\hspace{0.3cm}}lS[table-format=1.3]S[table-format=1.3]S[table-format=1.3]S[table-format=1.3]S[table-format=-1.3]@{\hspace{0.3cm}}}
        \toprule
        \textbf{Class} & {\textbf{Clean}} & {\textbf{FGSM $\epsilon=0.01$}} & {\textbf{FGSM $\epsilon=0.05$}} & {\textbf{FGSM $\epsilon=0.1$}} & {\textbf{$\Delta$ mAP ($\epsilon=0.1$)}} \\
        \midrule
        All            & 0.927 & 0.836 & 0.595 & 0.416 & -0.511 \\
        Green Light    & 0.927 & 0.783 & 0.617 & 0.459 & -0.468 \\
        Red Light      & 0.816 & 0.736 & 0.538 & 0.405 & -0.411 \\
        Speed Limit 10 & 0.863 & 0.584 & 0.000 & 0.000 & -0.863 \\
        Speed Limit 100 & 0.925 & 0.827 & 0.611 & 0.504 & -0.421 \\
        Speed Limit 110 & 0.852 & 0.821 & 0.555 & 0.000 & -0.852 \\
        Speed Limit 120 & 0.966 & 0.830 & 0.609 & 0.580 & -0.386 \\
        Speed Limit 20 & 0.991 & 0.946 & 0.791 & 0.543 & -0.448 \\
        Speed Limit 30 & 0.916 & 0.839 & 0.663 & 0.361 & -0.555 \\
        Speed Limit 40 & 0.978 & 0.934 & 0.662 & 0.508 & -0.470 \\
        Speed Limit 50 & 0.910 & 0.803 & 0.676 & 0.589 & -0.321 \\
        Speed Limit 60 & 0.951 & 0.875 & 0.312 & 0.245 & -0.706 \\
        Speed Limit 70 & 0.948 & 0.898 & 0.623 & 0.529 & -0.419 \\
        Speed Limit 80 & 0.975 & 0.855 & 0.706 & 0.426 & -0.549 \\
        Speed Limit 90 & 0.891 & 0.831 & 0.694 & 0.515 & -0.376 \\
        Stop           & 0.990 & 0.980 & 0.870 & 0.580 & -0.410 \\
        \bottomrule
    \end{tabular}
\end{table}


\subsubsection{PGD Attack on YOLO}

% Writing introduction paragraph
\paragraph{Introduction}\leavevmode\\
The Projected Gradient Descent (PGD) attack is an iterative adversarial attack that generates adversarial examples by repeatedly perturbing the input image along the gradient of the loss function, with projections to ensure perturbations remain within an $\ell_\infty$-bounded $\epsilon$-ball \cite{madry2018towardspgd}. For YOLO \cite{lu2017noneed}, PGD targets high-confidence detections from raw model outputs before non-maximum suppression (NMS), aiming to maximize classification errors and cause misclassifications or missed detections. Unlike the single-step FGSM, PGD's iterative nature with smaller step sizes and optional random initialization enhances attack strength while maintaining visual similarity, controlled by the perturbation magnitude $\epsilon$.

% Writing methodology paragraph with algorithm
\paragraph{Methodology}\leavevmode\\
\begin{breakablealgorithm}{Projected Gradient Descent (PGD) Attack on YOLO}
\begin{algorithmic}[1]
\Require $x \in \mathbb{R}^{H \times W \times C}$: input image, $model$: YOLO model, $\epsilon > 0$: perturbation magnitude, $T > 0$: number of iterations, $\alpha > 0$: step size, $\tau > 0$: confidence threshold, $random\_start$: boolean for random initialization
\Ensure $x_{adv} \in \mathbb{R}^{H \times W \times C}$: adversarial image
\State Preprocess $x$ for input to $model$ \Comment{E.g., normalize to $[0, 1]$, resize}
\State Initialize $x_{adv} \gets x$ \Comment{Set initial adversarial image}
\If{$random\_start$}
    \State Sample $\eta \sim \mathcal{U}(-\epsilon, \epsilon)^{H \times W \times C}$ \Comment{Random noise in $\epsilon$-ball}
    \State $x_{adv} \gets x + \eta$
    \State $x_{adv} \gets \text{clip}(x_{adv}, 0, 1)$ \Comment{Ensure pixel values in $[0, 1]$}
\EndIf
\For{$t = 1$ to $T$}
    \State Compute model predictions: $p \gets model(x_{adv})$ \Comment{$p$ includes bounding boxes, objectness scores, class probabilities}
    \State Extract detections $\mathcal{D} = \{ (b_i, s_i, c_i) \mid s_i \cdot \max(c_i) > \tau \}$, where $b_i$ are bounding boxes, $s_i$ are objectness scores, $c_i \in \mathbb{R}^K$ are class probabilities for $K$ classes
    \If{$\mathcal{D} = \emptyset$}
        \State \Return None \Comment{No high-confidence detections}
    \EndIf
    \State Select target class labels $y_i = \arg\max_k (c_i)_k$ for each $(b_i, s_i, c_i) \in \mathcal{D}$
    \State Define classification loss: $L(x_{adv}, \mathcal{D}) = \sum_{(b_i, s_i, c_i) \in \mathcal{D}} \ell(c_i, y_i)$, where $\ell$ is cross-entropy loss
    \State Compute gradient: $g \gets \nabla_{x_{adv}} L(x_{adv}, \mathcal{D}) \in \mathbb{R}^{H \times W \times C}$ \Comment{Gradient of loss w.r.t. $x_{adv}$}
    \State Update adversarial image: $x_{adv} \gets x_{adv} + \alpha \cdot \text{sign}(g)$
    \State Project to $\epsilon$-ball: $x_{adv} \gets \text{clip}(x_{adv}, x - \epsilon, x + \epsilon)$ \Comment{Ensure $||x_{adv} - x||_\infty \leq \epsilon$}
    \State Clamp to valid range: $x_{adv} \gets \text{clip}(x_{adv}, 0, 1)$ \Comment{Ensure pixel values in $[0, 1]$}
\EndFor
\State \Return $x_{adv}$
\end{algorithmic}
\end{breakablealgorithm}

% Writing results placeholder paragraph
\paragraph{Results}\leavevmode\\
The PGD attack was applied with perturbation magnitudes $\epsilon=0.01$, $0.05$, and $0.1$, using $num\_steps=10$ and $step\_size (\alpha)=epsilon/4$. Figure~\ref{fig:pgd_example_yolo} illustrates the effect of the attack on a sample image with $\epsilon=0.1$ and $\alpha = 0.25$, showing the original image, the adversarial image, and the pixel-wise difference.
\begin{figure}[H]
    \centering
    \includegraphics[width=0.8\textwidth]{images/CV/OD/pgd.png} % Replace with the actual path to your image file
    \caption{PGD Example showing the original, adversarial, and difference images.}
    \label{fig:pgd_example_yolo}
\end{figure}

  Although the relatively large epsilon ($\epsilon=0.1$), the adversarial image is not much different from the original image, as the perturbations are small, even smaller than ~\ref{fig:fgsm_example_yolo} and the original detections disappeared with the appearance of a new incorrect class. 

  
The impact of the PGD attack on YOLO's performance is detailed in Tables~\ref{tab:yolo_pgd_overall_performance} and \ref{tab:yolo_pgd_class_results}, which present general metrics (precision, recall, mAP@0.5, mAP@0.5:0.95) and class-specific mAP@0.5 across different $\epsilon$ values, highlighting the degradation in detection performance due to iterative adversarial perturbations.

% Table 1: Overall Performance
\begin{table}[H]
\caption{Overall Performance of YOLO under Clean and PGD Attacks}
    \label{tab:yolo_pgd_overall_performance}
    \centering
    \footnotesize
    \begin{tabular}{@{\hspace{0.3cm}}lS[table-format=1.3]S[table-format=1.3]S[table-format=1.3]S[table-format=1.3]@{\hspace{0.3cm}}}
        \toprule
        \textbf{Metric} & {\textbf{Clean}} & {\textbf{PGD $\epsilon=0.01$}} & {\textbf{PGD $\epsilon=0.05$}} & {\textbf{PGD $\epsilon=0.1$}} \\
        \midrule
        Precision     & 0.935 & 0.756 & 0.356 & 0.127 \\
        Recall        & 0.854 & 0.682 & 0.328 & 0.140 \\
        mAP@0.5       & 0.927 & 0.699 & 0.349 & 0.174 \\
        mAP@0.5:0.95  & 0.811 & 0.614 & 0.318 & 0.165 \\
        \bottomrule
    \end{tabular}
\end{table}

% Table 2: mAP@0.5 per Class
\begin{table}[H]
\caption{Class-Specific mAP@0.5 under Clean and PGD Attacks}
    \label{tab:yolo_pgd_class_results}
    \centering
    \footnotesize
    \begin{tabular}{@{\hspace{0.3cm}}lS[table-format=1.3]S[table-format=1.3]S[table-format=1.3]S[table-format=1.3]S[table-format=-1.3]@{\hspace{0.3cm}}}
        \toprule
        \textbf{Class} & {\textbf{Clean}} & {\textbf{PGD $\epsilon=0.01$}} & {\textbf{PGD $\epsilon=0.05$}} & {\textbf{PGD $\epsilon=0.1$}} & {\textbf{$\Delta$ mAP ($\epsilon=0.1$)}} \\
        \midrule
        Green Light     & 0.927 & 0.681 & 0.230 & 0.014 & -0.913 \\
        Red Light       & 0.816 & 0.639 & 0.294 & 0.098 & -0.718 \\
        Speed Limit 10  & 0.863 & 0.500 & 0.000 & 0.000 & -0.863 \\
        Speed Limit 100 & 0.925 & 0.874 & 0.369 & 0.128 & -0.797 \\
        Speed Limit 110 & 0.852 & 0.807 & 0.445 & 0.385 & -0.467 \\
        Speed Limit 120 & 0.966 & 0.842 & 0.521 & 0.387 & -0.579 \\
        Speed Limit 20  & 0.991 & 0.913 & 0.373 & 0.000 & -0.991 \\
        Speed Limit 30  & 0.916 & 0.850 & 0.286 & 0.163 & -0.753 \\
        Speed Limit 40  & 0.978 & 0.966 & 0.653 & 0.233 & -0.745 \\
        Speed Limit 50  & 0.910 & 0.812 & 0.448 & 0.409 & -0.501 \\
        Speed Limit 60  & 0.951 & 0.773 & 0.273 & 0.309 & -0.642 \\
        Speed Limit 70  & 0.948 & 0.854 & 0.492 & 0.338 & -0.610 \\
        Speed Limit 80  & 0.975 & 0.891 & 0.543 & 0.265 & -0.710 \\
        Speed Limit 90  & 0.891 & 0.773 & 0.418 & 0.407 & -0.484 \\
        Stop            & 0.990 & 0.918 & 0.547 & 0.269 & -0.721 \\
        \bottomrule
    \end{tabular}
\end{table}


\subsubsection{Dense Adversary Generation (DAG)}

\paragraph{Introduction}\leavevmode\\
The DAG algorithm \cite{xie2017DAG} was originally designed as a \textit{targeted} adversarial attack for object detection and semantic segmentation. In this work, we extend the framework to include both \textit{targeted} and \textit{untargeted} variants, implementing key modifications:

\noindent\textbf{Targeted Implementation:}
\begin{itemize}
    \item Follows the original formulation: Assigns specific adversarial labels $l'_n \neq l_n$ for each target
    \item Optimizes the loss $\mathcal{L} = \sum_{n} (f_{l_n}(X,t_n) - f_{l'_n}(X,t_n))$ to force misclassifications into predefined classes
    \item Maintains spatial perturbation structure through normalized gradient steps ($\gamma$)
\end{itemize}

\noindent\textbf{untargeted Implementation:}
\begin{itemize}
    \item Novel adaptation: Removes class-specific targeting
    \item Uses loss $\mathcal{L} = -\sum_{n} f_{l_n}(X,t_n)$ to degrade model confidence without directional guidance
    \item Requires modified convergence criteria as no specific target class exists
    \item Usually requires less conversion time as no target is specified.   
\end{itemize}

\paragraph{Methodology}\leavevmode\\

The following pseudocode outlines the unified DAG attack framework, with annotations to highlight the differences between the targeted and untargeted variants.
\begin{breakablealgorithm}{Dense Adversary Generation (DAG) Attack on YOLO}
\begin{algorithmic}[1]
\Require $x \in \mathbb{R}^{H \times W \times C}$: input image, $model$: YOLO model, $y' \in \{1, \dots, K\}$: adversarial class (targeted only), $T > 0$: number of iterations, $\gamma > 0$: step size, $\tau > 0$: confidence threshold
\Ensure $x_{adv} \in \mathbb{R}^{H \times W \times C}$: adversarial image, $success$ (targeted only): boolean success flag
\State Preprocess $x$ for input to $model$ \Comment{E.g., normalize to $[0, 1]$, resize}
\State Initialize $x_{adv} \gets x$ \Comment{Set initial adversarial image}
\State Compute initial predictions: $p \gets model(x_{adv})$ \Comment{$p$ includes bounding boxes, objectness scores, class probabilities}
\State Extract initial detections $\mathcal{D}_0 = \{ (b_i, s_i, c_i) \mid s_i \cdot \max(c_i) > \tau \}$, where $b_i$ are bounding boxes, $s_i \in [0, 1]$ are objectness scores, $c_i \in \mathbb{R}^K$ are class probabilities for $K$ classes
\State Define original classes: $Y_0 = \{ \arg\max_k (c_i)_k \mid (b_i, s_i, c_i) \in \mathcal{D}_0 \}$ \Comment{Targeted only}
\For{$t = 1$ to $T$}
    \State Compute predictions: $p \gets model(x_{adv})$ \Comment{Bounding boxes, objectness scores, class probabilities}
    \State Extract detections $\mathcal{D} = \{ (b_i, s_i, c_i) \mid s_i \cdot \max(c_i) > \tau \}$
    \If{$\mathcal{D} = \emptyset$}
        \If{targeted attack}
            \State Define loss: $L(x_{adv}) = \sum_{i} c_i[y']$ \Comment{Boost adversarial class globally}
        \Else
            \State \textbf{break} \Comment{Stop if no detections (untargeted)}
        \EndIf
    \Else
        \If{targeted attack}
            \State Define loss: $L(x_{adv}, \mathcal{D}) = \sum_{(b_i, s_i, c_i) \in \mathcal{D}} (c_i[y'] - c_i[\arg\max_k (c_i)_k])$ \Comment{Encourage adversarial class}
        \Else
            \State Define loss: $L(x_{adv}, \mathcal{D}) = -\sum_{(b_i, s_i, c_i) \in \mathcal{D}} c_i[\arg\max_k (c_i)_k]$ \Comment{Reduce confidence in predicted classes}
        \EndIf
    \EndIf
    \State Compute gradient: $g \gets \nabla_{x_{adv}} L \in \mathbb{R}^{H \times W \times C}$ \Comment{Gradient of loss w.r.t. $x_{adv}$}
    \State Update adversarial image: $x_{adv} \gets x_{adv} + \gamma \cdot \text{sign}(g)$
    \State Clamp to valid range: $x_{adv} \gets \text{clip}(x_{adv}, 0, 1)$ \Comment{Ensure pixel values in $[0, 1]$}
    \If{targeted attack}
        \State Extract current detections $\mathcal{D}' = \{ (b_i, s_i, c_i) \mid s_i \cdot \max(c_i) > \tau \}$
        \State Define current classes: $Y' = \{ \arg\max_k (c_i)_k \mid (b_i, s_i, c_i) \in \mathcal{D}' \}$
        \If{$y' \in Y'$ and $Y' \cap Y_0 = \emptyset$}
            \State \textbf{break} \Comment{Stop if adversarial class achieved and original classes removed}
        \EndIf
    \EndIf
\EndFor
\If{targeted attack}
    \State Extract final detections $\mathcal{D}_f = \{ (b_i, s_i, c_i) \mid s_i \cdot \max(c_i) > \tau \}$
    \State Define final classes: $Y_f = \{ \arg\max_k (c_i)_k \mid (b_i, s_i, c_i) \in \mathcal{D}_f \}$
    \State $success \gets (y' \in Y_f) \land (Y_f \cap Y_0 = \emptyset)$
    \State \Return $x_{adv}, success$
\Else
    \State \Return $x_{adv}$
\endIf
\end{algorithmic}
\end{breakablealgorithm}

\paragraph{Results}
The DAG attack was applied with $\gamma=0.003$ unlike the original paper\cite{xie2017DAG} that used $\gamma=0.5$ that caused too much distortion when tested in YOLO, and $num\_iterations=20$ (a sufficient number of iterations is defined by trying) . Figures~\ref{fig:TDAG_example_yolo} and ~\ref{fig:UDAG_example_yolo} illustrate the effect of the attack (targeted and non-targeted) on a sample image with $\gamma=0.003$, $num\_iterations=20$ and green light as target for the targeted variant. showing the original image, the adversarial images, and the pixel-wise differences.
\begin{figure}[H]
    \centering
    \includegraphics[width=0.8\textwidth]{images/CV/OD/TDAG.png} % Replace with the actual path to your image file
    \caption{Targeted DAG Example showing the original, adversarial, and difference images.}
    \label{fig:TDAG_example_yolo}
\end{figure}


\begin{figure}[H]
    \centering
    \includegraphics[width=0.8\textwidth]{images/CV/OD/UDAG.png} % Replace with the actual path to your image file
    \caption{Untargeted DAG Example showing the original, adversarial, and difference images.}
    \label{fig:UDAG_example_yolo}
\end{figure}

Both targeted and untargeted attacks show visually small perturbations~\ref{fig:TDAG_example_yolo}~\ref{fig:UDAG_example_yolo}. The targeted attack successfully suppressed the original detection and boosted the target one, while the untargeted attack did its job and suppressed the original detection. 

  
The Dense Adversary Generation (DAG) attack both targeted and untargeted overall performance are summarized in Tables ~\ref{tab:yolo_dag_untargeted_overall_performance} and ~\ref{tab:yolo_dag_targeted_overall_performance} and the Class-Specific mAP@0.5 results are in Tables ~\ref{tab:yolo_dag_untargeted_class_results} and ~\ref{tab:yolo_dag_targeted_class_results}


% Table 1: Untargeted Overall Performance
\begin{table}[H]
    \centering
    \caption{Overall Performance of YOLO under Clean and Untargeted DAG Attacks}
    \label{tab:yolo_dag_untargeted_overall_performance}
    \footnotesize
    \begin{tabular}{@{\hspace{0.3cm}}lS[table-format=1.3]S[table-format=1.3]S[table-format=1.3]S[table-format=1.3]S[table-format=1.3]S[table-format=1.3]@{\hspace{0.3cm}}}
        \toprule
        \textbf{Metric} & {\textbf{Clean}} & {\textbf{Untargeted DAG}} \\
        \midrule
        Precision & 0.935 & 0.693 \\
        Recall & 0.854 & 0.385 \\
        mAP@0.5 & 0.927 & 0.547 \\
        mAP@0.5:0.95 & 0.811 & 0.441 \\
        \bottomrule
    \end{tabular}
\end{table}

% Table 2: Untargeted Class-Specific mAP@0.5
\begin{table}[H]
    \centering
    \caption{Class-Specific mAP@0.5 under Clean and Untargeted DAG Attacks}
    \label{tab:yolo_dag_untargeted_class_results}
    \footnotesize
    \begin{tabular}{@{\hspace{0.3cm}}lS[table-format=1.3]S[table-format=1.3]S[table-format=-1.3]@{\hspace{0.3cm}}}
        \toprule
        \textbf{Class} & {\textbf{Clean}} & {\textbf{Untargeted DAG}} & {\textbf{$\Delta$ mAP}} \\
        \midrule
        Green Light    & 0.927 & 0.325 & -0.602 \\
        Red Light      & 0.816 & 0.480 & -0.336 \\
        Speed Limit 10 & 0.863 & 0.667 & -0.196 \\
        Speed Limit 20 & 0.920 & 0.516 & -0.404 \\
        Speed Limit 30 & 0.916 & 0.533 & -0.383 \\
        Speed Limit 40 & 0.925 & 0.614 & -0.311 \\
        Speed Limit 50 & 0.966 & 0.710 & -0.256 \\
        Speed Limit 60 & 0.910 & 0.440 & -0.470 \\
        Speed Limit 70 & 0.860 & 0.590 & -0.270 \\
        Speed Limit 80 & 0.978 & 0.500 & -0.478 \\
        Speed Limit 90 & 0.961 & 0.616 & -0.345 \\
        Speed Limit 100 & 0.925 & 0.538 & -0.387 \\
        Speed Limit 110 & 0.952 & 0.416 & -0.536 \\
        Speed Limit 120 & 0.966 & 0.543 & -0.423 \\
        Stop           & 0.990 & 0.807 & -0.183 \\
        \bottomrule
    \end{tabular}
\end{table}

% Table 3: Targeted Overall Performance
\begin{table}[H]
    \centering
    \caption{Overall Performance of YOLO under Clean and Targeted DAG Attacks}
    \label{tab:yolo_dag_targeted_overall_performance}
    \footnotesize
    \begin{tabular}{@{\hspace{0.3cm}}lS[table-format=1.3]S[table-format=1.3]S[table-format=1.3]S[table-format=1.3]S[table-format=1.3]S[table-format=1.3]@{\hspace{0.3cm}}}
        \toprule
        \textbf{Metric} & {\textbf{Clean}} & {\textbf{Targeted DAG}} \\
        \midrule
        Precision & 0.935 & 0.148 \\
        Recall & 0.854 & 0.0931 \\
        mAP@0.5 & 0.927 & 0.0947 \\
        mAP@0.5:0.95 & 0.811 & 0.0819 \\
        \bottomrule
    \end{tabular}
\end{table}

% Table 4: Targeted Class-Specific mAP@0.5
\begin{table}[H]
    \centering
    \caption{Class-Specific mAP@0.5 under Clean and Targeted DAG Attacks}
    \label{tab:yolo_dag_targeted_class_results}
    \footnotesize
    \begin{tabular}{@{\hspace{0.3cm}}lS[table-format=1.3]S[table-format=1.3]S[table-format=-1.3]@{\hspace{0.3cm}}}
        \toprule
        \textbf{Class} & {\textbf{Clean}} & {\textbf{Targeted DAG}} & {\textbf{$\Delta$ mAP}} \\
        \midrule
        Green Light    & 0.927 & 0.0125 & -0.9145 \\
        Red Light      & 0.863 & 0.0659 & -0.7971 \\
        Speed Limit 10 & 0.863 & 0.0679 & -0.7951 \\
        Speed Limit 100 & 0.852 & 0.0679 & -0.7841 \\
        Speed Limit 110 & 0.921 & 0.321 & -0.600 \\
        Speed Limit 120 & 0.953 & 0.0733 & -0.8797 \\
        Speed Limit 20 & 0.96 & 0.0733 & -0.8867 \\
        Speed Limit 30 & 0.95 & 0.0737 & -0.8763 \\
        Speed Limit 40 & 0.943 & 0.0755 & -0.8675 \\
        Speed Limit 50 & 0.95 & 0.0737 & -0.8763 \\
        Speed Limit 60 & 0.943 & 0.0763 & -0.8667 \\
        Speed Limit 70 & 0.943 & 0.0771 & -0.8659 \\
        Speed Limit 80 & 0.943 & 0.0801 & -0.8629 \\
        Speed Limit 90 & 0.989 & 0.0801 & -0.9089 \\
        Stop           & 0.989 & 0.234 & -0.755 \\
        \bottomrule
    \end{tabular}
\end{table}

\subsubsection{"Adversarial Examples that Fool Detectors" Attack}
\paragraph{Introduction}\leavevmode\\
The "Adversarial Examples that Fool Detectors" method, introduced by Lu et al. \cite{lu2017fooldetectors}, presents a technique for creating adversarial examples designed specifically to deceive object detection models.The authors demonstrate its effectiveness by generating robust digital and physical perturbations for objects like stop signs and faces, showing successful generalization from Faster R-CNN to YOLO.

The core of the attack is an iterative optimization process that seeks to minimize the mean confidence score of all detection proposals corresponding to a specific target class. To ensure the adversarial perturbation remains visually subtle, the objective function includes an L2 regularization term that penalizes deviations from the original image. The image is updated at each step using the sign of the gradient (an L-infinity projected gradient descent step) to minimize this objective.

In this work, we implement the attack strategy detailed in \cite{lu2017fooldetectors}. We adapt their stop sign-focused methodology to target traffic sign classes within the LISA (Laboratory for Intelligent \& Safe Automobiles) dataset. Our goal is to evaluate the effectiveness of this attack on a fine-tuned YOLOv8 model.
\paragraph{Methodology}\leavevmode\\
The attack is implemented as an iterative process targeting specific traffic sign classes within an input image, aiming to reduce the YOLOv8 model's confidence in detecting those classes.
The algorithm is formalized as follows:


\begin{breakablealgorithm}{Adversarial Examples that Fool Detectors Attack on YOLO}
\begin{algorithmic}[1]
\Require $x \in \mathbb{R}^{H \times W \times C}$: input image, $model$: YOLO model, $Y_t \subseteq \{1, \dots, K\}$: target classes to suppress, $T > 0$: number of iterations, $\gamma > 0$: step size, $\tau > 0$: confidence threshold, $\lambda_{reg} > 0$: L2 regularization weight
\Ensure $x_{adv} \in \mathbb{R}^{H \times W \times C}$: adversarial image
\State Preprocess $x$ for input to $model$ \Comment{E.g., normalize to $[0, 1]$, resize}
\State Initialize $x_{adv} \gets x$ \Comment{Set initial adversarial image}
\State $x_0 \gets x$ \Comment{Save original image}
\For{$t \gets 1$ to $T$}
    \State Compute predictions: $p \gets model(x_{adv})$ \Comment{$p$ includes bounding boxes, objectness scores, class probabilities}
    \State Extract detections $\mathcal{D} = \{ (b_i, s_i, c_i) \mid s_i \cdot \max(c_i) > \tau \}$, where $b_i$ are bounding boxes, $s_i \in [0, 1]$ are objectness scores, $c_i \in \mathbb{R}^K$ are class probabilities for $K$ classes
    \State Define detected classes: $Y_d = \{ \arg\max_k (c_i)_k \mid (b_i, s_i, c_i) \in \mathcal{D} \}$
    \If{$Y_t \cap Y_d = \emptyset$}
        \State \textbf{break} \Comment{Stop if no target classes detected}
    \EndIf
    \State Extract target confidences: $c_{i,t} = c_i[y]$ for each $(b_i, s_i, c_i) \in \mathcal{D}$, $y \in Y_t$
    \State Define detection loss: $L_{det}(x_{adv}, \mathcal{D}) = \frac{1}{|\mathcal{D}| \cdot |Y_t|} \sum_{(b_i, s_i, c_i) \in \mathcal{D}} \sum_{y \in Y_t} c_{i,t}$ \Comment{Mean confidence of target classes}
    \State Define L2 penalty: $L_{reg}(x_{adv}) = \lambda_{reg} \cdot ||x_{adv} - x_0||_2^2$ \Comment{L2 distance from original image}
    \State Define total loss: $L(x_{adv}, \mathcal{D}) = L_{det}(x_{adv}, \mathcal{D}) + L_{reg}(x_{adv})$
    \State Compute gradient: $g \gets \nabla_{x_{adv}} L(x_{adv}, \mathcal{D}) \in \mathbb{R}^{H \times W \times C}$ \Comment{Gradient of loss w.r.t. $x_{adv}$}
    \State Update adversarial image: $x_{adv} \gets x_{adv} - \gamma \cdot \text{sign}(g)$ \Comment{L-infinity PGD step}
    \State Clamp to valid range: $x_{adv} \gets \text{clip}(x_{adv}, 0, 1)$ \Comment{Ensure pixel values in $[0, 1]$}
\EndFor
\State \Return $x_{adv}$ \Comment{Return adversarial image}
\end{algorithmic}
\end{breakablealgorithm}

\paragraph{Results}\leavevmode\\
The "Adversarial Examples that Fool Detectors" attack was applied with $\gamma=0.01$ and $\lambda = 0.01$ within the ranges suggested by the original paper \cite{lu2017fooldetectors}, with $num\_iterations=10$ (a sufficient number of iterations is defined by trying). Figure~\ref{fig:fool_example_yolo} illustrates the effect of the attack (targeted and non-targeted) on a sample image with $\gamma=0.01$, $\lambda = 0.01$ and $num\_iterations=10$, showing the original image, the adversarial images, and the pixel-wise differences.
\begin{figure}[H]
    \centering
    \includegraphics[width=0.8\textwidth]{images/CV/OD/fool.png} % Replace with the actual path to your image file
    \caption{"Adversarial Examples that Fool Detectors" Example showing the original, adversarial, and difference images.}
    \label{fig:fool_example_yolo}
\end{figure}

The adversarial image is not much different from the original image, as the perturbations are visually small, and the original detections disappeared with the appearance of a new incorrect class.

The impact of the "Adversarial Examples that Fool Detectors" attack on YOLO's performance is detailed in Tables~\ref{tab:yolo_fool_overall_performance} and \ref{tab:yolo_fool_class_results}, which present general metrics (precision, recall, mAP@0.5, mAP@0.5:0.95) and class-specific mAP@0.5, highlighting the degradation in detection performance due to iterative adversarial perturbations.

\begin{table}[H]
    \centering
    \caption{Overall Performance of YOLO under Clean and "Adversarial Examples that Fool Detectors" Attack}
    \label{tab:yolo_fool_overall_performance}
    \footnotesize
    \begin{tabular}{@{\hspace{0.3cm}}lS[table-format=1.3]S[table-format=1.3]S[table-format=1.3]S[table-format=1.3]S[table-format=1.3]S[table-format=1.3]@{\hspace{0.3cm}}}
        \toprule
        \textbf{Metric} & {\textbf{Clean}} & {\textbf{Adversarial Attack}} \\
        \midrule
        Precision & 0.935 & 0.238 \\
        Recall & 0.854 & 0.16 \\
        mAP@0.5 & 0.927 & 0.155 \\
        mAP@0.5:0.95 & 0.811 & 0.125 \\
        \bottomrule
    \end{tabular}
\end{table}

\begin{table}[H]
    \centering
    \caption{Class-Specific mAP@0.5 under Clean and "Adversarial Examples that Fool Detectors" Attack}
    \label{tab:yolo_fool_class_results}
    \footnotesize
    \begin{tabular}{@{\hspace{0.3cm}}lS[table-format=1.3]S[table-format=1.3]S[table-format=-1.3]@{\hspace{0.3cm}}}
        \toprule
        \textbf{Class} & {\textbf{Clean}} & {\textbf{Adversarial Attack}} & {\textbf{$\Delta$ mAP}} \\
        \midrule
        All            & 0.927 & 0.155 & -0.772 \\
        Green Light    & 0.927 & 0.0204 & -0.906 \\
        Red Light      & 0.816 & 0.213 & -0.603 \\
        Speed Limit 10 & 0.863 & 0.000 & -0.863 \\
        Speed Limit 20 & 0.913 & 0.160 & -0.753 \\
        Speed Limit 30 & 0.937 & 0.0471 & -0.890 \\
        Speed Limit 40 & 0.966 & 0.220 & -0.746 \\
        Speed Limit 50 & 0.916 & 0.136 & -0.780 \\
        Speed Limit 60 & 0.953 & 0.0842 & -0.869 \\
        Speed Limit 70 & 0.866 & 0.132 & -0.734 \\
        Speed Limit 80 & 0.947 & 0.246 & -0.701 \\
        Speed Limit 90 & 0.806 & 0.229 & -0.577 \\
        Speed Limit 100 & 0.941 & 0.0695 & -0.871 \\
        Speed Limit 110 & 0.959 & 0.223 & -0.736 \\
        Speed Limit 120 & 0.966 & 0.125 & -0.841 \\
        Stop           & 0.990 & 0.421 & -0.569 \\
        \bottomrule
    \end{tabular}
\end{table}


\subsubsection{Universal Adversarial Perturbations (UAP)}
\paragraph{Introduction}\leavevmode\\
Universal adversarial perturbations (UAP) paper shows existence of universal image-agnostic perturbations for state-of-the-art deep neural networks \cite{moosavi2017universal}. The algorithm searches a universal perturbation for a set of training points. In this work, we modified the original algorithm to be more suitable to find universal adversarial perturbations that can fool object detectors.

\paragraph{Methodology}\leavevmode\\
The methodology adapts the universal adversarial perturbation approach from Moosavi-Dezfooli et al. \cite{moosavi2017universal} to target object detection models, specifically using a YOLOv8 model. The process trains a single perturbation vector, termed delta, that, when added to any input image, reduces the detection confidence across multiple images. The key steps involve iterative optimization over a set of training images, incorporating a regularization term to maintain perturbation subtlety. The algorithm was trained with an epsilon of 0.1 and 10 epochs to constrain the perturbation magnitude and utilized a dataset with 15 classes. The training dataset was constructed by sampling 100 images per class, ensuring a balanced representation. The algorithm is formalized as follows:

\begin{breakablealgorithm}{Universal Adversarial Perturbation Training for YOLO}
\begin{algorithmic}[1]
\Require $X = \{x_1, \dots, x_N\} \subset \mathbb{R}^{H \times W \times C}$: training images, $model$: YOLO model, $E > 0$: number of epochs, $\lambda_{reg} > 0$: L2 regularization weight, $\epsilon > 0$: perturbation bound, $\tau > 0$: confidence threshold, $\alpha > 0$: learning rate
\Ensure $\delta \in \mathbb{R}^{H \times W \times C}$: universal perturbation
\State Preprocess each $x_i \in X$ for input to $model$ \Comment{E.g., normalize to $[0, 1]$, resize}
\State Initialize $\delta \gets 0$ \Comment{Set initial perturbation to zero}
\For{$e \gets 1$ to $E$}
    \State Initialize epoch loss: $L_e \gets 0$
    \For{$x_i \in X$}
        \State Compute adversarial image: $x_{adv} \gets x_i + \delta$
        \State Compute predictions: $p \gets model(x_{adv})$ \Comment{$p$ includes bounding boxes, objectness scores, class probabilities}
        \State Extract detections $\mathcal{D} = \{ (b_j, s_j, c_j) \mid s_j \cdot \max(c_j) > \tau \}$, where $b_j$ are bounding boxes, $s_j \in [0, 1]$ are objectness scores, $c_j \in \mathbb{R}^K$ are class probabilities for $K$ classes
        \If{$\mathcal{D} = \emptyset$}
            \State \textbf{continue} \Comment{Skip if no high-confidence detections}
        \EndIf
        \State Define target classes: $y_j = \arg\max_k (c_j)_k$ for each $(b_j, s_j, c_j) \in \mathcal{D}$
        \State Define classification loss: $L_{cls}(x_{adv}, \mathcal{D}) = \sum_{(b_j, s_j, c_j) \in \mathcal{D}} \ell(c_j, y_j)$, where $\ell$ is cross-entropy loss
        \State Define regularization loss: $L_{reg}(\delta) = \lambda_{reg} \cdot ||\delta||_2^2$ \Comment{L2 norm of perturbation}
        \State Define total loss: $L(x_{adv}, \delta, \mathcal{D}) = L_{cls}(x_{adv}, \mathcal{D}) + L_{reg}(\delta)$
        \State Compute gradient: $g \gets \nabla_{\delta} L(x_{adv}, \delta, \mathcal{D}) \in \mathbb{R}^{H \times W \times C}$ \Comment{Gradient w.r.t. $\delta$}
        \State Update perturbation: $\delta \gets \delta - \alpha \cdot g$ \Comment{SGD update}
        \State Clamp perturbation: $\delta \gets \text{clip}(\delta, -\epsilon, \epsilon)$ \Comment{Ensure $||\delta||_\infty \leq \epsilon$}
        \State Accumulate epoch loss: $L_e \gets L_e + L(x_{adv}, \delta, \mathcal{D})$
    \EndFor
    \State Compute average epoch loss: $L_e \gets L_e / N$
\EndFor
\State \Return $\delta$
\end{algorithmic}
\end{breakablealgorithm}


\paragraph{Results}\leavevmode\\
Figure~\ref{fig:uap_example_yolo} illustrates the effect of the attack on a sample image showing the original image, the adversarial images, and the pixel-wise differences.
\begin{figure}[H]
    \centering
    \includegraphics[width=0.8\textwidth]{images/CV/OD/UAP_1.png} % Replace with the actual path to your image file
    \caption{UAP Example showing the original, adversarial, and difference images.}
    \label{fig:uap_example_yolo}
\end{figure}

The adversarial image is not much different from the original image, as the perturbations are visually small, and the original detections disappeared.

The impact of the UAP attack on YOLO's performance is detailed in Tables~\ref{tab:yolo_uap_overall_performance} and \ref{tab:yolo_uap_class_results}, which present general metrics (precision, recall, mAP@0.5, mAP@0.5:0.95) and class-specific mAP@0.5, highlighting the degradation in detection performance due to iterative adversarial perturbations.

\begin{table}[H]
    \centering
    \caption{Overall Performance of YOLO under Clean and UAP Attack}
    \label{tab:yolo_uap_overall_performance}
    \footnotesize
    \begin{tabular}{@{\hspace{0.3cm}}lS[table-format=1.3]S[table-format=1.3]S[table-format=1.3]S[table-format=1.3]S[table-format=1.3]S[table-format=1.3]@{\hspace{0.3cm}}}
        \toprule
        \textbf{Metric} & {\textbf{Clean}} & {\textbf{UAP Attack}} \\
        \midrule
        Precision & 0.935 & 0.661 \\
        Recall & 0.854 & 0.197 \\
        mAP@0.5 & 0.927 & 0.443 \\
        mAP@0.5:0.95 & 0.811 & 0.365 \\
        \bottomrule
    \end{tabular}
\end{table}

\begin{table}[H]
    \centering
    \caption{Class-Specific mAP@0.5 under Clean and UAP Attack}
    \label{tab:yolo_uap_class_results}
    \footnotesize
    \begin{tabular}{@{\hspace{0.3cm}}lS[table-format=1.3]S[table-format=1.3]S[table-format=-1.3]@{\hspace{0.3cm}}}
        \toprule
        \textbf{Class} & {\textbf{Clean}} & {\textbf{UAP Attack}} & {\textbf{$\Delta$ mAP}} \\
        \midrule
        all            & 0.927 & 0.443 & -0.484 \\
        Green Light    & 0.927 & 0.454 & -0.473 \\
        Red Light      & 0.816 & 0.533 & -0.283 \\
        Speed Limit 10 & 0.863 & 0.000 & -0.863 \\
        Speed Limit 20 & 0.913 & 0.265 & -0.648 \\
        Speed Limit 30 & 0.937 & 0.526 & -0.411 \\
        Speed Limit 40 & 0.966 & 0.543 & -0.423 \\
        Speed Limit 50 & 0.916 & 0.526 & -0.390 \\
        Speed Limit 60 & 0.953 & 0.416 & -0.537 \\
        Speed Limit 70 & 0.866 & 0.487 & -0.379 \\
        Speed Limit 80 & 0.947 & 0.578 & -0.369 \\
        Speed Limit 90 & 0.806 & 0.603 & -0.203 \\
        Speed Limit 100 & 0.941 & 0.506 & -0.435 \\
        Speed Limit 110 & 0.959 & 0.000 & -0.959 \\
        Speed Limit 120 & 0.966 & 0.558 & -0.408 \\
        Stop           & 0.990 & 0.570 & -0.420 \\
        \bottomrule
    \end{tabular}
\end{table}


\subsection{Black Box Attacks}
\subsubsection{Square Attack}
\paragraph{Introduction}\leavevmode\\
The Square Attack, proposed by Andriushchenko et al.~\cite{andriushchenko2020square}, is a score-based, black-box adversarial attack that crafts $\ell_\infty$- and $\ell_2$-norm-bounded perturbations via random search. By applying localized, square-shaped updates, it achieves high query efficiency and success rates, outperforming state-of-the-art black-box attacks by up to a factor of three in query consumption on image classification tasks~\cite{andriushchenko2020square}. Its gradient-independent nature makes it robust against defenses like gradient masking, ideal for testing real-world vision systems.

In this work, we adapt the Square Attack to target the YOLOv8 object detection model on the LISA traffic sign dataset. Our implementation employs $\ell_\infty$-norm perturbations with $\epsilon \in \{0.05, 0.1, 0.2\}$ and patch sizes of \{256, 512\}, using a cross-entropy loss to suppress detection confidences below a threshold of \SI{0.5}{}. Evaluated via metrics like mAP50, the attack significantly degrades performance from a baseline of \SI{0.927}{} to as low as \SI{0.597}{}.


% Defining the methodology for Square Attack implementation
\paragraph{Methodology}\leavevmode\\
We implement the Square Attack~\cite{andriushchenko2020square} aiming to suppress detection confidences below a threshold, using a cross-entropy loss on class logits of detected objects. For each image, we initialize the adversarial image as the original and iterate up to N times. In each iteration, a random square patch (size 256 or 512) is selected, and a uniform perturbation $\delta \in [-\epsilon, \epsilon]$ ($\epsilon \in \{\SI{0.05}{},\SI{0.1}{}, \SI{0.2}{}\}$) is applied, followed by clamping to the $\ell_\infty$-ball and $[0,1]$. Updates are accepted if they increase the loss, with early stopping if no detections exceed the threshold. Unlike the original attack’s image classification focus, dynamic square sizes, and margin loss~\cite{andriushchenko2020square}, our adaptation uses fixed patch sizes for LISA’s high-resolution images, a cross-entropy loss for object detection, and no specific initialization. We evaluate performance via mAP50, mAP50-95, precision, and recall, comparing baseline (mAP50: \SI{0.927}{}) to adversarial results.

% Defining the Square Attack algorithm for YOLOv8
\begin{breakablealgorithm}{Square Attack for Object Detection}
\label{alg:square_attack}
\begin{algorithmic}[1]
\State \textbf{Input:} Image $\mathbf{x} \in [0,1]^{3 \times H \times W}$, model $f$, norm bound $\epsilon$, patch size $s$, max iterations $N$, confidence threshold $\tau$
\State \textbf{Output:} Adversarial image $\hat{\mathbf{x}}$
\State $\hat{\mathbf{x}} \gets \mathbf{x}$ 
\State $L_{\text{best}} \gets 0$ \Comment{Track best cross-entropy loss}
\For{$i = 1$ to $N$}
    \State Sample patch position $(x, y)$ 
    \State Sample $\delta \sim \text{Uniform}[-\epsilon, \epsilon]^{3 \times s \times s}$
    \State $\hat{\mathbf{x}}_{\text{new}} \gets \hat{\mathbf{x}}$
    \State Apply $\delta$ to patch at $(x, y)$ in $\hat{\mathbf{x}}_{\text{new}}$
    \State $\hat{\mathbf{x}}_{\text{new}} \gets \text{Clamp}(\hat{\mathbf{x}}_{\text{new}}, \mathbf{x} - \epsilon, \mathbf{x} + \epsilon)$ \Comment{Ensure $\ell_\infty$ bound}
    \State $\hat{\mathbf{x}}_{\text{new}} \gets \text{Clamp}(\hat{\mathbf{x}}_{\text{new}}, 0, 1)$
    \State Compute detections $D = f(\hat{\mathbf{x}}_{\text{new}})$
    \State Compute loss $L = \sum_{\text{det} \in D} \text{CrossEntropy}(\text{det}_{\text{logits}}, \text{det}_{\text{class}})$
    \If{$L > L_{\text{best}}$} \Comment{Increase loss vs. decrease in~\cite{andriushchenko2020square}}
        \State $\hat{\mathbf{x}} \gets \hat{\mathbf{x}}_{\text{new}}$
        \State $L_{\text{best}} \gets L$
    \EndIf
    \If{No detection in $D$ has confidence $> \tau$} \Comment{Early stopping for object detection}
        \State \textbf{break}
    \EndIf
\EndFor
\State \textbf{return} $\hat{\mathbf{x}}$
\end{algorithmic}
\end{breakablealgorithm}


\paragraph{Results}\leavevmode\\
Table~\ref{tab:yolo_square_overall_performance} presents the performance of YOLOv8 on the LISA dataset under clean and Square Attack conditions with $\epsilon \in \{\SI{0.05}{}, \SI{0.1}{}, \SI{0.2}{}\}$ and patch sizes \{256, 512\}. The attack significantly degrades performance, with mAP50 dropping to \SI{0.597}{} and Recall to \SI{0.404}{} in the worst case ($\epsilon=\SI{0.2}{}$, patch size=512), a \SI{35.6}{\%} relative mAP50 decrease. Larger $\epsilon$ and patch sizes amplify degradation, though patch size 256 outperforms 512 at $\epsilon=\SI{0.1}{}$ (mAP50: \SI{0.821}{} vs. \SI{0.795}{}). Compared to Andriushchenko et al.~\cite{andriushchenko2020square}, who report near-\SI{100}{\%} misclassification success on ImageNet at $\epsilon=\SI{0.05}{}$ using dynamic patch sizes, our object detection task shows milder degradation (mAP50: \SI{0.889}{}--\SI{0.890}{}) due to multiple objects per image. Our fixed patch sizes to deal with LISA dataset and cross-entropy loss, with \SI{500}{} iterations, achieve comparable query efficiency to the paper’s \SI{500}{}--\SI{1500}{} queries. These results highlight YOLOv8’s vulnerability to $\ell_\infty$-norm perturbations.

The Attacks parameters are (epsilon, patch size)

\begin{table}[H]
\centering
\caption{Detection Performance for Each (epsilon, patch size) Configuration}
\label{tab:yolo_square_overall_performance}
\begin{tabular}{l c c c c c c c}
\toprule
\textbf{Metric} 
    & Clean & (0.05, 256) & (0.05, 512) & (0.1, 256) & (0.1, 512) & (0.2, 256) & (0.2, 512) \\
\midrule
\textbf{Precision} & 0.935 & 0.919 & 0.858 & 0.801 & 0.824 & 0.761 & 0.728 \\
\textbf{Recall}    & 0.854 & 0.812 & 0.848 & 0.759 & 0.697 & 0.536 & 0.404 \\
\textbf{mAP50}     & 0.927 & 0.889 & 0.890 & 0.821 & 0.795 & 0.681 & 0.597 \\
\textbf{mAP50-95}  & 0.811 & 0.776 & 0.771 & 0.708 & 0.674 & 0.579 & 0.499 \\
\bottomrule
\end{tabular}
\end{table}


\subsection{Physical Attacks}

\section{Defenses}
\subsection{Adversarial Training}
As a foundational defense, adversarial training is highly effective for improving the robustness of object detection models. have been thoroughly explained in Section~\ref{sec:adversarial_training} of our chapter on Image Classification (Chapter~\ref{chap:classification}).

In the context of object detection, the same core idea applies. Adversarial examples are generated to maximize the detector's composite loss function, which includes errors in both class prediction and bounding box localization. By training the model on these worst-case examples, it learns to produce more stable bounding boxes and maintain correct class labels, even when subjected to carefully crafted perturbations.

\subsubsection{Training Data Composition}
To build a model with comprehensive robustness against a wide array of potential threats, we adopted a mixed-data adversarial training strategy. Rather than training solely against a single attack like PGD, we augmented our original dataset with adversarial examples generated from a diverse suite of attack algorithms.

Our baseline training set consists of 10,590 clean images from the LISA dataset. We expanded this set by generating 8,252 adversarial examples, creating a final combined training corpus of 18,842 images. The composition of the adversarial portion was intentionally heterogeneous to encourage the model to learn more generalized defensive features. The breakdown is as follows:

\begin{itemize}
    \item \textbf{Projected Gradient Descent (PGD):} 2,000 images were generated using a strong, iterative PGD attack to build foundational robustness.
    
    \item \textbf{Fast Gradient Sign Method (FGSM):} 3,000 images were created using the fast, single-step FGSM attack. Including these ensures the model is resilient against less computationally expensive threats.
    
    \item \textbf{Targeted DAG (DAG):} 500 images were perturbed with the DAG attack, which focuses on creating misclassifications to specific target labels.
    
    \item \textbf{Universal Adversarial Perturbations (UAP):} 1,500 examples were generated using a single, image-agnostic perturbation, testing the model's resilience to systematic noise patterns.
    
    \item \textbf{FoolDetectors:} 1,500 examples were generated using an attack specifically designed to fool object detectors by causing bounding boxes to vanish or mis-localize.
\end{itemize}

This diverse mixture of gradient-based, targeted, universal, and detector-specific attacks ensures that the model does not overfit its defense to one particular threat vector. The entire combined dataset of clean and adversarial images was then shuffled and used for training the YOLOv8 model, ensuring that each batch contained a random combination of samples and forced the model to learn robust features effective against multiple attack modalities.
\subsection{Dynamic Defensive Preprocessing Pipeline}
\label{sec:defense_pipeline}

In addition to hardening the model itself via adversarial training, a complementary approach is to defend the system at the input level. Instead of processing every image as a potential threat, we designed a more efficient, two-stage dynamic pipeline. This pipeline first uses a lightweight classifier to detect likely adversarial examples and then applies a series of purification transformations only to those flagged images. This targeted approach aims to neutralize threats while minimizing the computational overhead and potential quality degradation for benign images.

The entire process is illustrated in Figure~\ref{fig:defense_pipeline_flowchart} and the specific parameters used are summarized in Table~\ref{tab:pipeline_params}.
\begin{figure}[H]
    \centering
    \includegraphics[width=0.3\textwidth]{images/CV/OD/flow.pdf} % You would create this image
    \caption{The logic of our two-stage dynamic defense pipeline. An input image is first classified. If it is deemed adversarial, it undergoes purification before being sent to the YOLOv8 detector. Otherwise, it is passed through directly.}
    \label{fig:defense_pipeline_flowchart}
\end{figure}
% It's highly recommended to create a simple flowchart for this.
% You can use tools like draw.io and export it as a PDF/PNG.

\subsubsection{Stage 1: Adversarial Detection}

% Introducing the adversarial detection stage
The initial stage of our pipeline serves as a critical gatekeeper to identify and mitigate adversarial perturbations in input images. We utilize a \texttt{MobileNetV2} model, configured as a binary classifier, to distinguish between clean and adversarially perturbed images. The \texttt{MobileNetV2} architecture was selected for its computational efficiency and low inference latency, enabling rapid preprocessing without introducing significant bottlenecks in the pipeline.

% Describing the classification process
Each incoming image is evaluated by the \texttt{MobileNetV2} classifier. If the predicted probability of the image being adversarial exceeds a predefined threshold of 0.3, the image is flagged for subsequent purification. Conversely, images classified as clean are directly forwarded to the YOLOv8 detector for object detection without modification.

\paragraph{Classifier Training Process}

% Outlining the dataset composition
\subparagraph{Dataset Composition}
The training dataset comprises 14,000 images, including 7,000 clean images and 7,000 adversarially perturbed images. The adversarial subset consists of 2,000 images generated using Projected Gradient Descent (PGD), 2,000 using Fast Gradient Sign Method (FGSM), 1,500 using Universal Adversarial Perturbations (UAP), 500 using Targeted Deep Adversarial Generation (DAG), and 1,000 images from food detection datasets. The dataset is split into 12,000 images for training, 1,000 for validation, and 1,000 for testing.

% Specifying training parameters
\subparagraph{Training Parameters}
The \texttt{MobileNetV2} model was trained for 10 epochs to achieve robust performance in distinguishing clean and adversarial images.

\paragraph{Classifier Performance}

% Presenting test dataset performance
\subparagraph{Test Dataset Performance}
On the test dataset, the classifier achieved an accuracy of 0.9980. demonstrating high reliability in controlled settings.

% Evaluating performance on real-world data
\subparagraph{Real-World Data Performance}
To assess generalization, the classifier was evaluated on a separate dataset of 1,100 images, comprising 500 clean images and 600 adversarial images generated using attack methods not encountered during training. With the adversarial classification threshold set at 0.3, the classifier correctly identified 555 out of 600 adversarial images, yielding an accuracy of 92.5\%.


\subsubsection{Stage 2: Input Purification Transformations}
When an image is flagged as adversarial, it is subjected to a sequence of three transformations designed to disrupt or remove the adversarial noise:

\begin{enumerate}
    \item \textbf{JPEG Compression:} The image is first compressed using the JPEG algorithm. This lossy compression method is particularly effective at removing high-frequency noise, which is a common characteristic of adversarial perturbations.
    
    \item \textbf{Gaussian Blurring:} Following compression, a Gaussian blur is applied. This acts as a low-pass filter, smoothing the image and further mitigating any fine-grained, pixel-level manipulations that may have survived the compression step.
    
    \item \textbf{Random Resizing:} Finally, the image is randomly resized by a scale factor drawn uniformly from a predefined range. This transformation disrupts the precise spatial structure of an adversarial pattern, as the attack was crafted for the original image dimensions.
\end{enumerate}


These transformations, applied in sequence, create a powerful defense by attacking the adversarial perturbation from multiple angles: frequency, local pixel coherence, and spatial alignment.

\begin{table}[h!]
    \centering
    \caption{Hyperparameters for the Input Purification Stage}
    \label{tab:pipeline_params}
    \begin{tabular}{ll}
        \toprule
        \textbf{Parameter} & \textbf{Value} \\
        \midrule
        JPEG Quality & 30 \\
        Gaussian Blur Sigma ($\sigma$) & 1.5 \\
        Random Resize Scale Range & (0.8, 1.2) \\
        Adversarial Classifier Threshold & 0.3 \\
        \bottomrule
    \end{tabular}
\end{table}

\subsubsection{Maintaining Label Integrity: Bounding Box Rescaling}
A critical implementation detail of this pipeline is ensuring the ground truth remains accurate. Since the resizing transformation alters the image dimensions, the original YOLO bounding box labels become invalid. 

To correct this, we implemented a label rescaling function. After an image is resized, this function mathematically adjusts the normalized coordinates (\texttt{x\_center}, \texttt{y\_center}, \texttt{width}, \texttt{height}) of each bounding box to correspond to the new image dimensions. This ensures that the YOLOv8 model is evaluated against the correct ground truth for the transformed image, guaranteeing a fair assessment of the defense's effectiveness.


\subsubsection{Performance Overhead and Latency}

While effectiveness against attacks is the primary goal, a practical defense must also be computationally efficient to be viable in real-world applications. To quantify the overhead of our dynamic defensive pipeline, we measured the processing time on our test set. The key performance metrics are summarized in Table~\ref{tab:pipeline_timing}.

\begin{table}[h!]
    \centering
    \caption{Per-Image Timing Analysis of the Dynamic Defense Pipeline}
    \label{tab:pipeline_timing}
    \begin{tabular}{l S[table-format=3.1]} % Using siunitx for decimal alignment
        \toprule
        \textbf{Metric} & {\textbf{Time (ms)}} \\
        \midrule
        \multicolumn{2}{l}{\textit{Cost for All Images:}} \\
        \quad Classification Overhead & 21.2 \\
        \quad Label & 0.2 \\
        \addlinespace % Adds a little space
        \multicolumn{2}{l}{\textit{Cost for Flagged Adversarial Images Only:}} \\
        \quad Purification (JPEG, Blur, Resize) & 204.6 \\
        \addlinespace
        \midrule
        \textbf{Total Latency (per Clean Image)} & \textbf{21.7} \\
        \textbf{Total Latency (per Adversarial Image)} & \textbf{226.0} \\
        \bottomrule
    \end{tabular}
\end{table}

The results clearly illustrate the primary advantage of our dynamic approach. Benign images, which constitute the vast majority of inputs in a real-world scenario, incur a minimal overhead of only 21.7 ms, primarily from the initial classification check. This ensures that the system's performance on clean data is not significantly degraded.

In contrast, when an image is flagged as adversarial, the full purification pipeline is triggered, resulting in a total processing time of 226.0 ms. While this latency is substantially higher, it is only applied to a small subset of suspicious inputs. This demonstrates a successful trade-off, concentrating the computational effort where it is most needed to neutralize threats while maintaining high throughput for normal operations. The total time to process the entire test set was approximately 2 minutes and 18 seconds.

\ifSubfilesClassLoaded{
  \bibliographystyle{unsrt}
  \bibliography{references}
}{}



\end{document}
